% !TeX root = ../thuthesis-example.tex

\chapter{引言}

\section{项目目的}

随着移动互联网应用的发展，用户每天接触到的消息、新闻、群组动态和社交内容不断增加。传统按时间倒序展示内容的方式实现简单，但在内容数量较大时容易出现两个问题：一是用户真正关心的内容被大量低相关内容淹没；二是系统无法根据用户兴趣、关注关系和近期行为动态调整内容排序。推荐系统的引入可以在“内容生产者”和“内容消费者”之间建立更有效的匹配关系，使用户在较短时间内看到更有价值的内容。

本项目以类Telegram社交应用为业务背景，设计一套基于Spring Boot的分布式推荐系统。项目目标不是单纯完成一个接口，而是从课程设计角度形成完整的工程闭环：前端负责页面展示和交互，Spring Boot后端负责接口、业务编排、推荐算法与安全控制，Supabase PostgreSQL负责云端持久化数据，Redis负责缓存和时间线加速，Apifox负责接口联调和测试记录。

本项目重点解决以下问题：

\begin{enumerate}
  \item 构建符合分布式应用课程要求的后端架构。后端采用Spring Boot进行模块划分，按控制层、服务层、仓储层和算法策略层组织代码，避免把接口、业务逻辑和数据访问混在同一个文件中。
  \item 建立真实可解释的推荐算法流程。推荐系统不是只按浏览量排序，而是综合用户关注关系、内容标签、行为事件、兴趣向量、热点内容、冷启动策略和实验配置进行候选生成与排序。
  \item 接入云端数据库。系统使用Supabase PostgreSQL保存用户、群组、联系人、新闻内容、用户行为、兴趣向量和实验配置，使数据结构具备远程访问、持久化和后续扩展能力。
  \item 完成接口联调与测试。通过Apifox构造推荐Feed接口请求，验证请求参数、接口路径和返回结构，为后续自动化测试和接口文档共享打基础。
\end{enumerate}

从工程角度看，本项目将推荐逻辑拆分为多个阶段，每个阶段只承担清晰的职责：召回阶段负责扩大候选集，过滤阶段负责保证候选安全和有效，评分阶段负责计算相关性，选择阶段负责控制数量和多样性，副作用阶段负责记录曝光与实验数据。这种分层流水线思想有利于后续替换算法、扩展数据源和定位问题。

\section{本文组织结构}

本文按照课程大作业报告模板组织内容，章节安排如下：

第一章为引言，说明项目背景、项目目的和报告组织结构。

第二章为相关技术概述，介绍Spring Boot、Spring Cloud思想、Supabase PostgreSQL、Redis、Apifox、React前端以及个性化推荐算法等关键技术。

第三章为系统分析设计，围绕需求分析、可行性分析、系统架构、模块设计、数据库设计和推荐算法设计展开说明，并给出系统架构图、推荐算法流程图和数据库实体关系图。

第四章为项目测试，说明接口测试环境、Apifox联调过程、核心接口测试方案、数据库检查结果以及测试中发现的安全风险。

第五章为总结与展望，总结本项目完成的主要工作，分析目前仍然存在的不足，并提出后续在微服务拆分、模型训练、安全策略和监控运维方面的改进方向。
